\documentclass[10.5pt,a4paper]{article}

\usepackage[margin=0.75in]{geometry}
\usepackage{mathptmx}
\usepackage{titlesec}
\usepackage{enumitem}
\usepackage{xcolor}
\usepackage[hidelinks]{hyperref}
\usepackage{tabularx}

\definecolor{accent}{HTML}{8C1D18}

\pagestyle{empty}

\titleformat{\section}{\large\bfseries\color{accent}}{}{0em}{}[\vspace{-0.3em}{\color{accent}\hrule height 0.8pt}\vspace{0.4em}]
\titlespacing{\section}{0pt}{0.8em}{0.3em}

\setlist[itemize]{leftmargin=1.1em, itemsep=1pt, topsep=1pt, parsep=0pt}

\newcommand{\heading}[2]{%
  \noindent\textbf{#1} \hfill #2\par
}

\begin{document}

% ---------- HEADER ----------
\begin{center}
  {\Huge \bfseries Rafat Abdullah}\\[0.4em]
  Dhaka, Bangladesh \;\textbar\; 01715138888 \;\textbar\; \href{mailto:rafatabdullah4425@gmail.com}{rafatabdullah4425@gmail.com}\\[0.2em]
  \href{https://www.linkedin.com/in/rafat-pantho/}{linkedin.com/in/rafat-pantho} \;\textbar\; \href{https://github.com/Rafat-Pantho}{github.com/Rafat-Pantho}
\end{center}

% ---------- OBJECTIVE ----------
\section*{Objective}
Computer Science undergraduate seeking an internship at Huawei Technologies Bangladesh to apply hands-on
experience in software development, applied AI, and systems-level programming to real-world engineering
problems, while growing within a large-scale technology organization.

% ---------- EDUCATION ----------
\section*{Education}
\heading{Islamic University of Technology (IUT), Gazipur, Bangladesh}{Expected Graduation: 2027}
B.Sc. in Computer Science and Engineering (CSE) \hfill Currently in 3rd Year, 2nd Semester\\
CGPA: 2.75 / 4.00

% ---------- SKILLS ----------
\section*{Technical Skills}
\textbf{Programming Languages:} Python, C, C++\\
\textbf{Database Compatibility:} MySQL, Oracle, MongoDB, PostgreSQL

% ---------- PROJECTS ----------
\section*{Projects}

\heading{VisionCart --- Visual Inventory Management System}{\href{https://github.com/Rafat-Pantho/VisionCart}{GitHub}}
\textit{Python, FastAPI, SQLite/SQLAlchemy, Google GenAI SDK (Gemma vision-language model)}
\begin{itemize}
  \item Built a full-stack system that replaces manual shelf counting: a photo of a shelf is analyzed and
        reconciled against a live inventory database in real time.
  \item Integrated a vision-language model (Gemma) via Google's GenAI SDK to detect products and count units
        directly from images, automatically flagging low-stock and out-of-stock items.
  \item Designed a FastAPI backend with a SQLAlchemy-backed database and a responsive vanilla JS/HTML/CSS
        dashboard (live KPIs, search/filter/sort, light \& dark themes); deployed live on Render.
\end{itemize}

\heading{My Pocket Buddy --- On-Device AI Companion App}{\href{https://github.com/Rafat-Pantho/my-pocket-buddy}{GitHub}}
\textit{TypeScript, React Native, Expo, llama.rn (on-device LLM inference), MMKV}
\begin{itemize}
  \item Built a cross-platform (iOS/Android) mobile app that runs a local LLM entirely on-device via
        \texttt{llama.rn} (GGUF models, CPU-only inference) for private, offline AI chat with no cloud API
        dependency.
  \item Designed a hybrid context-injection system for the chat engine: a sliding window of recent turns
        plus regex-triggered full-history recall, so the AI companion can reference earlier conversations.
  \item Implemented a voice-activated alarm that listens via microphone to detect when the user wakes up and
        dismiss it automatically, alongside a stats dashboard and an AI-generated exam-practice screen.
\end{itemize}

\heading{PMVC-WNM --- Banglish Sentiment Classification (Co-training Prototype)}{\href{https://github.com/Rafat-Pantho/ML-Banglish-co-training-prototype}{GitHub}}
\textit{Python, scikit-learn --- Course project, CSE 4622 Machine Learning Lab, IUT}
\begin{itemize}
  \item Implemented a two-view co-training pipeline (character n-gram SVM + phonetic-encoding logistic
        regression) to classify sentiment in Banglish (Romanized Bangla-English code-mixed) text as a
        lightweight alternative to deep learning.
  \item Designed a rule-based phonetic encoder to normalize inconsistent Romanized Bangla spellings (e.g.
        \textit{khacchi} / \textit{khacci} / \textit{khachi}) into shared phonetic codes.
  \item Added a weakly supervised noise model that down-weights unreliable pseudo-labels using cross-view
        disagreement, aiming for competitive accuracy with no GPU and orders of magnitude fewer parameters
        than a deep-learning baseline.
\end{itemize}

\heading{Local Judge for Windows --- Offline Competitive Programming Judge}{\href{https://github.com/Rafat-Pantho/localjudge-for-windows}{GitHub}}
\textit{Python, C++, JavaScript}
\begin{itemize}
  \item Built a local judge that compiles and runs Python/C++ submissions against test cases under
        configurable time and memory limits, without depending on any Linux-only tooling.
  \item Implemented real-time process time/memory tracking and verdict reporting (Accepted, Wrong Answer,
        Runtime Error, Time Limit Exceeded), mirroring the behavior of an online judge for local, offline
        practice on Windows.
\end{itemize}

% ---------- ACHIEVEMENTS ----------
\section*{Achievements}
\begin{itemize}
  \item \textbf{Competitive Programming:} Codeforces rating 885 (handle: \href{https://codeforces.com/profile/NebulaKnight}{NebulaKnight}); selected for the NHSPC national round twice
  \item \textbf{Hackathons:} IUT-RS Hackathon; Build with Gemma --- ML, AI, Deep Learning \& NLP Community (Kaggle)
\end{itemize}

\end{document}
